\chapter{conclusão}

O uso de dados de saúde em aplicações de machine learning apresenta desafios significativos relacionados à fragmentação das fontes, heterogeneidade semântica, viés nos dados e requisitos rigorosos de segurança e privacidade. No contexto do diagnóstico dermatológico, essas limitações afetam diretamente a eficácia dos modelos preditivos, sobretudo em populações sub-representadas. A coleta e integração sistemática de dados clínicos e imagens de lesões, respeitando princípios de governança e diversidade, torna-se essencial para o desenvolvimento de algoritmos mais justos, confiáveis e clinicamente relevantes. Assim, este trabalho parte do desafio de construir uma infraestrutura capaz de orquestrar esses dados de forma ética e interoperável, viabilizando avanços em inteligência artificial aplicada à saúde pública.

O desenvolvimento da plataforma DermAlert representa um avanço significativo na aplicação de arquitetura Data Fabric para o contexto dermatológico brasileiro, pois combina uma arquitetura resiliente e adaptável às condições de infraestrutura desafiadoras do sistema de saúde, uma camada de integração semântica que resolve a heterogeneidade terminológica entre instituições e uma demonstração prática de como tecnologias open-source podem ser ajustadas para criar soluções especializadas de baixo custo voltadas à saúde pública. Além de disponibilizar a infraestrutura, a plataforma possui um aplicativo para a aquisição de imagens de lesões de pele e metadados da lesão do paciente. Essa nova base de dados, respeitando os padrões internacionais e disponibilizada no Data Fabric, contribui com a comunidade científica com uma diversidade de tons de pele, possibilitando o treinamento de modelos de machine learning mais adaptados para a realidade da população brasileira.
Trabalhos futuros incluem expansão para outras subespecialidades dermatológicas além de lesões oncológicas, como dermatoses inflamatórias e infecciosas. Estamos desenvolvendo capacidade de análise longitudinal para acompanhamento automático de evolução de lesões ao longo do tempo e integrando modelos de risco personalizado baseados em fatores genéticos e ambientais.

Uma oportunidade particularmente promissora é a implementação de federação de aprendizado, permitindo que modelos sejam treinados em dados distribuídos sem necessidade de centralização, endereçando importantes preocupações de privacidade.

\section{Evoluções Futuras}
Com base na arquitetura completa proposta, algumas funcionalidades planejadas para uma futura versão da plataforma incluem o desenvolvimento de filtros personalizados para facilitar a criação de novos datasets com base em critérios específicos definidos pelo usuário. Também está planejada a ampliação do suporte a diferentes tipos de conexões com novas fontes de dados, aumentando a flexibilidade e a abrangência da plataforma. Outro avanço será a implementação de um modelo de permissões baseado em RBAC (Role-Based Access Control), permitindo um controle mais refinado de acesso aos dados e datasets. Por fim, o suporte ao mapeamento e relacionamento entre tabelas de uma mesma fonte, possibilitando análises mais completas e estruturadas dentro do próprio contexto da origem dos dados.
